\subsection{Expressions}
\label{Section 1.1.1}

One easy way to get started at programming is to examine some typical
interactions with an interpreter for the Python language. Imagine that you
are sitting at a computer terminal. You type an \newterm{expression}, and
the interpreter responds by displaying the result of its \newterm{evaluating}
that expression.

One kind of primitive expression you might type is a number. (More precisely,
the expression that you type consists of the numerals that represent the number
in base 10.) If you present Python with a number

\begin{codeBlock}
>>> 486
~\textit{486}~
\end{codeBlock}

\noindent
the interpreter will respond by printing \textit{486}\footnote{Throughout this
book, when we wish to emphasize the distinction between the input typed by the
user and the response printed by the interpreter, we will show the latter in
slanted characters.}

Expressions representing numbers may be combined with an expression
representing a primitive procedure (such as \code{+} or \code{*}) to form a
compound expression that represents the application of the procedure to those
numbers. For example:

\begin{codeBlock}
>>> 137 + 349
~\textit{486}~
\end{codeBlock}

\begin{codeBlock}
>>> 1000 - 334
~\textit{666}~
\end{codeBlock}

\begin{codeBlock}
>>> 5 * 99
~\textit{495}~
\end{codeBlock}

\begin{codeBlock}
>>> 10 / 5
~\textit{2.0}~
\end{codeBlock}

\begin{codeBlock}
>>> 2.7 + 10
~\textit{12.7}~
\end{codeBlock}

\noindent
Expressions such as these, formed by joining simpler expressions with a symbol
that denotes a procedure, are called \newterm{compound expressions}. The symbol
is called the \newterm{operator}, and the expressions it stands between are
called the \newterm{operands}. The value of a compound expression is obtained by
applying the procedure specified by the operator to the \newterm{arguments} that
are the values of the operands.\footnote{Python reserves the phrase
\newterm{compound statement} for \code{if}, \code{for}, \code{while} and their
kin. A compound statement directs the order in which things are done; a compound
expression denotes a value. We will have little to say about statements until we
need them.}

The convention of placing the operator in between its operands is known as
\newterm{infix notation}, and it is the way most common programming languages
work.  Infix is not the only way to represent the interaction between
operators and operands.  Lisp famously uses something called \newterm{prefix
notation}, for example \code{(- 1000 334)}. A few things are worth noting here:
in Lisp every operator is applied to its operands in the same way, so \code{-}
and \code{minus} would behave identically if defined identically; and by putting
the operator first and enclosing the whole expression in parentheses, no
ambiguity about the order of operations can arise. Infix, by contrast, needs a
rule for that order.  Python tries to closely model the
arithmetic rules we learn in school, PEMDAS\footnote{Parentheses, Exponents,
Multiplication and Division, and finally Addition and Subtraction. However,
there is one wrinkle, if you write \pycode{-2 ** 2} you get \pycode{-4} because
Python interprets it as \pycode{-(2 ** 2)} not \pycode{(-2)**2}.}

The original text notes that one of the advantages of prefix notation is that it
can accommodate procedures that take an arbitrary number of operands. Infix
cannot do this.  To add four numbers we have to do the following:

\begin{codeBlock}
>>> 21 + 35 + 12 + 7
~\textit{75}~
\end{codeBlock}

\noindent
And to multiply three numbers we need to do:

\begin{codeBlock}
>>> 25 * 4 * 12
~\textit{1200}~
\end{codeBlock}

\noindent
One advantage of infix notation is that expressions read much as they do in
ordinary arithmetic, for example:

\begin{codeBlock}
>>> 3 * 5 + 10 - 6
~\textit{19}~
\end{codeBlock}

\noindent
Contrast that with the original text's:

\begin{codeBlock}
(+ (* 3 5) (- 10 6))
\end{codeBlock}

\noindent
That is arguably much harder to read. Note that we are always free to group
expressions with parentheses in Python; the above could be written as:

\begin{codeBlock}
>>> (3 * 5) + (10 - 6)
~\textit{19}~
\end{codeBlock}

\noindent
Expressions can be as complicated as we want, consider

\begin{codeBlock}
>>> (((2 * 4) + (3 + 5)) * 3) + ((10 - 7) + 6)
~\textit{57}~
\end{codeBlock}

\noindent
which is hard enough to read in infix notation; it is even harder in prefix. The
original text notes that Lisp programmers follow a formatting convention called
\newterm{pretty-printing} to make such expressions easier to read. In Python the
convention is not merely discouraged but mechanically undone: the formatters the
community has settled on will collapse such an expression back onto one
line.\footnote{Given the expression above, laid out over five aligned lines,
\code{ruff format} returns it as the single line printed here. Where an
expression is too long to fit, the formatter breaks it at the loosest operator
and indents every part equally, so the indentation never records the shape of
the expression the way pretty-printing does.}

Even with complex expressions, the interpreter always operates in the same
basic cycle: It reads an expression from the terminal, evaluates the
expression, and prints the result.  This mode of operation is often expressed
by saying that the interpreter runs in a \newterm{read-eval-print loop}.
Observe in particular that it is not necessary to explicitly instruct the
interpreter to print the value of the expression.\footnote{Python mostly obeys
the convention that every expression has a value. That convention, which it
shares with Lisp, together with the old reputation of Lisp as an inefficient
language, is the source of the quip by Alan Perlis (paraphrasing Oscar Wilde)
that ``Lisp
programmers know the value of everything but the cost of nothing.'' It might be
said of Python programmers as well.}
